\documentclass[12pt,a4paper]{article}

\usepackage[T2A]{fontenc}
\usepackage[utf8]{inputenc}
\usepackage[russian]{babel}
\usepackage{amsmath,amssymb,amsthm,mathtools}
\usepackage{geometry}
\usepackage{microtype}
\usepackage{enumitem}
\usepackage{xcolor}
\usepackage{tikz}
\usepackage{hyperref}

\geometry{left=25mm,right=20mm,top=22mm,bottom=24mm}
\hypersetup{
  colorlinks=true,
  linkcolor=blue!45!black,
  urlcolor=blue!55!black
}
\setlist{nosep}
\setlength{\parindent}{1.25em}
\setlength{\parskip}{0.15em}

\newtheorem{theorem}{Теорема}
\newtheorem{proposition}{Утверждение}
\newtheorem{lemma}{Лемма}
\theoremstyle{definition}
\newtheorem{definition}{Определение}
\theoremstyle{remark}
\newtheorem*{remark}{Замечание}

\DeclareMathOperator{\Sp}{Sp}
\DeclareMathOperator{\diag}{diag}
\DeclareMathOperator{\cond}{cond}
\newcommand{\CC}{\mathbb{C}}
\newcommand{\RR}{\mathbb{R}}
\newcommand{\ZZ}{\mathbb{Z}}
\newcommand{\eps}{\varepsilon}
\newcommand{\norm}[1]{\left\lVert #1\right\rVert}

\DeclareMathOperator{\ind}{ind}
\newcommand{\Conv}{\mathcal C}

\title{Критерий обратимости\\оператора Винера--Хопфа}
\author{Конспект лекции}
\date{7 ноября 2026 года}

\begin{document}
\maketitle

\section{Сжатые сдвиги на положительной полуоси}

Сохраним обозначения предыдущих лекций:
\[
  X=\ell^1(\ZZ),\qquad
  (P_+x)_k=\begin{cases}x_k,&k\geqslant1,\\0,&k\leqslant0,\end{cases}
  \qquad X_+=\operatorname{Im}P_+.
\]
Подпространство $X_+$ отождествляем с $\ell^1(\mathbb N)$,
где $\mathbb N=\{1,2,\ldots\}$. Двусторонний сдвиг
$V$ задан формулой $(Vx)_k=x_{k-1}$; его символ равен $t$.
Для $\kappa\in\ZZ$ рассмотрим
\[
  S_\kappa=(P_+V^\kappa P_+)\big|_{X_+}.
\]
Все утверждения об обратимости сжатий относятся к $X_+$.
На этом пространстве $P_+$ действует как единичный оператор.

Обозначим односторонние сдвиги
\[
  R(x_1,x_2,\ldots)=(0,x_1,x_2,\ldots),\qquad
  L(x_1,x_2,\ldots)=(x_2,x_3,\ldots).
\]
Тогда $S_1=R$, $S_{-1}=L$, $LR=I$, а $RL=I-E_1$,
где $E_1$ оставляет только первую координату.

\begin{theorem}
Тип обратимости $S_\kappa$ определяется знаком $\kappa$:
\[
  \begin{array}{c|c|c}
    \kappa& S_\kappa&\text{обратимость}\\\hline
    0&I&\text{двусторонняя}\\
    \kappa>0&R^\kappa&\text{только слева}\\
    \kappa<0&L^{-\kappa}&\text{только справа}
  \end{array}
\]
\end{theorem}

\begin{proof}
При $\kappa=0$ имеем $S_0=I$ на $X_+$.

Пусть $\kappa=m>0$. Тогда
\[
  S_mx=(\underbrace{0,\ldots,0}_{m},x_1,x_2,\ldots),\qquad
  S_{-m}S_m=L^mR^m=I.
\]
Значит, $S_{-m}$ является левым обратным к $S_m$.
Однако
\[
  \operatorname{Im}S_m
  =\{y\in X_+:y_1=\cdots=y_m=0\}\ne X_+.
\]
Оператор не сюръективен, поэтому правого обратного не имеет.

Пусть теперь $\kappa=-m<0$. Тогда
\[
  S_{-m}x=(x_{m+1},x_{m+2},\ldots),\qquad S_{-m}S_m=I.
\]
Таким образом, $S_m$ --- правый обратный к $S_{-m}$.
Но $S_{-m}e^{(1)}=0$, где $e^{(1)}\ne0$. Оператор не инъективен
и потому не имеет левого обратного.
\end{proof}

Более точно, из координатных формул следует
\[
  \ker S_{-m}=\operatorname{span}\{e^{(1)},\ldots,e^{(m)}\},\qquad
  S_mS_{-m}=I-E_m,
\]
где $E_m$ оставляет первые $m$ координат. Эти равенства уточняют
причины отсутствия второго обратного.

\section{Достаточные условия: ненулевой символ и его индекс}

Пусть $a\in W$, где $W$ --- алгебра Винера,
\[
  a(t)=\sum_{k\in\ZZ}a_kt^k,\qquad
  \norm{a}_W=\sum_{k\in\ZZ}|a_k|<\infty.
\]
Оператор свёртки $A$ имеет ядро $(a_k)$, а соответствующий оператор
Винера--Хопфа равен
\[
  T_a=(P_+AP_+)\big|_{X_+},\qquad
  (T_ax)_k=\sum_{j=1}^{\infty}a_{k-j}x_j\quad(k\geqslant1).
\]

\begin{theorem}
Пусть $a(t)\ne0$ для всех $t\in\Gamma=\{t:|t|=1\}$ и
$\kappa=\ind a$. Тогда
\begin{enumerate}
  \item при $\kappa>0$ оператор $T_a$ обратим только слева;
  \item при $\kappa=0$ оператор $T_a$ двусторонне обратим;
  \item при $\kappa<0$ оператор $T_a$ обратим только справа.
\end{enumerate}
\end{theorem}

\begin{proof}
По факторизации предыдущей лекции
\[
  T_a=T_-S_\kappa T_+,
\]
где $T_-$ и $T_+$ двусторонне обратимы на $X_+$.
Умножение слева или справа на обратимый оператор сохраняет тип
односторонней обратимости. Поэтому тип $T_a$ совпадает с типом
$S_\kappa$, уже установленным в первом разделе.
\end{proof}

Для пояснения можно явно выписать один обратный с нужной стороны.
При $\kappa=m>0$ левым обратным является
\[
  C=T_+^{-1}S_{-m}T_-^{-1},\qquad CT_a=I.
\]
При $\kappa=-m<0$ правым обратным является
\[
  D=T_+^{-1}S_mT_-^{-1},\qquad T_aD=I.
\]
При $\kappa=0$ получаем $T_a^{-1}=T_+^{-1}T_-^{-1}$.
Эти формулы разворачивают применение факторизации.

\begin{remark}
Пока доказан результат только для ненулевого символа.
Чтобы получить критерий, необходимо исключить одностороннюю
обратимость при наличии нуля на окружности. Этому посвящена
оставшаяся часть лекции.
\end{remark}

\section{Два линейных символа с нулём на окружности}

Понадобится устойчивость обратимости в операторной норме:
множества двусторонне обратимых, обратимых только слева и обратимых
только справа операторов открыты. Первое утверждение следует
из ряда Неймана, два других доказаны на предыдущей лекции.

\subsection{Символ \texorpdfstring{$\alpha_0-t$}{alpha0 - t}}

\begin{lemma}
Если $|\alpha_0|=1$, то оператор
\[
  T_{\alpha_0-t}=\alpha_0 I-R
\]
не имеет ни левого, ни правого обратного.
\end{lemma}

\begin{proof}
Рассмотрим семейство $F_\alpha=T_{\alpha-t}=\alpha I-R$.
Разность операторов равна скалярному множителю единицы:
\[
  \norm{F_\alpha-F_{\alpha_0}}=|\alpha-\alpha_0|.
\]
Если $|\alpha|<1$, у многочлена $\alpha-t$ один нуль внутри круга,
поэтому $\ind(\alpha-t)=1$. Значит, $F_\alpha$ обратим только слева.
Если $|\alpha|>1$, индекс равен нулю и $F_\alpha$ двусторонне обратим.

В любой окрестности $\alpha_0$ существуют точки обоих типов,
например $(1-\varepsilon)\alpha_0$ и $(1+\varepsilon)\alpha_0$.
Если бы $F_{\alpha_0}$ был двусторонне обратим, то все достаточно
близкие операторы были бы двусторонне обратимы, что противоречит
поведению при $|\alpha|<1$.
Если бы он был обратим только слева или только справа, то все
достаточно близкие операторы имели бы тот же тип. Это противоречит
двусторонней обратимости при $|\alpha|>1$.
Все возможные виды односторонней обратимости исключены.
\end{proof}

\begin{figure}[ht]
  \centering
  \begin{tikzpicture}[>=stealth,scale=1.8]
    \draw[->,gray!60] (-1.25,0) -- (1.7,0);
    \draw[->,gray!60] (0,-1.2) -- (0,1.3);
    \draw[blue!60!black,thick] (0,0) circle (1);
    \draw[dashed] (1,0) circle (0.35);
    \fill (1,0) circle (0.025) node[above right] {$\alpha_0$};
    \fill (0.77,0) circle (0.025) node[below left] {$\alpha_{\rm in}$};
    \fill (1.23,0) circle (0.025) node[below right] {$\alpha_{\rm out}$};
    \node[above left] at (0,0) {$0$};
    \node[above] at (0,1) {$|\alpha|=1$};
  \end{tikzpicture}
  \caption{Любая окрестность точки окружности содержит внутренние и внешние точки}
\end{figure}

В рукописи приведена идея возмущения параметра. Здесь отдельно
проверены все три возможных типа обратимости, чтобы вывод
<<ни с одной стороны>> был полностью обоснован.

\subsection{Символ \texorpdfstring{$1-\alpha_0/t$}{1 - alpha0/t}}

\begin{lemma}
Если $|\alpha_0|=1$, то оператор
\[
  T_{1-\alpha_0/t}=I-\alpha_0 L
\]
не имеет ни левого, ни правого обратного.
\end{lemma}

\begin{proof}
Положим $G_\alpha=I-\alpha L$. Поскольку $\norm{L}=1$,
\[
  \norm{G_\alpha-G_{\alpha_0}}=|\alpha-\alpha_0|.
\]
При $|\alpha|<1$ функция $1-\alpha/t$ не обращается в нуль на
окружности и имеет индекс $0$; оператор $G_\alpha$ двусторонне обратим.
Это также прямо следует из ряда Неймана.
При $|\alpha|>1$ запишем
\[
  1-\frac\alpha t=\frac{t-\alpha}{t}.
\]
Нулей внутри круга нет, а полюс в нуле один, поэтому индекс равен $-1$.
Оператор $G_\alpha$ обратим только справа.

Если бы $G_{\alpha_0}$ был двусторонне обратим, получили бы противоречие
с близкими внешними точками. Если бы он был обратим только с одной
стороны, получили бы противоречие с близкими внутренними точками,
где оператор двусторонне обратим. Устойчивость каждого типа
обратимости завершает доказательство.
\end{proof}

Доказательство второй леммы в рукописи не развёрнуто;
здесь оно приведено по аналогии с первой.

\section{Необходимое условие для конечного ядра}

\begin{definition}
В терминологии рукописи оператор свёртки называется \emph{финитным},
если его ядро имеет конечное число ненулевых коэффициентов.
Тогда символ является многочленом Лорана:
\[
  a(t)=\sum_{k=-m}^{n}a_kt^k=\frac{p(t)}{t^m},\qquad m,n\geqslant0.
\]
\end{definition}

Матрица такой свёртки ленточная, но по-прежнему бесконечная.
Конечность ядра не означает конечного ранга оператора:
например, ядро с единственным коэффициентом $a_0=1$ задаёт
тождественный оператор на бесконечномерном пространстве.

\begin{theorem}
Если ядро $A$ конечно и $T_a$ обратим хотя бы с одной стороны,
то $a(t)\ne0$ для всех $t\in\Gamma$.
\end{theorem}

\begin{proof}
Предположим, что $a(\alpha_0)=0$ и $|\alpha_0|=1$.
Если $a\equiv0$, оператор нулевой и не имеет односторонних обратных,
поэтому этот случай сразу исключён. В остальных случаях
\[
  p(t)=(t-\alpha_0)q(t),\qquad
  a(t)=\frac{(t-\alpha_0)q(t)}{t^m}.
\]
Рассмотрим отдельно левую и правую обратимость.

\emph{1. Левый обратный.} Выделим множитель
\[
  a(t)=b(t)(\alpha_0-t),\qquad b(t)=-\frac{q(t)}{t^m}.
\]
Оба множителя имеют конечные ядра; символ $\alpha_0-t$ принадлежит
$W_+$, то есть соответствующий оператор $A_0$ нижнетреуголен.
По лемме о проекторе $A_0P_+=P_+A_0P_+$, поэтому
\[
  T_a=T_bT_{\alpha_0-t}.
\]
Если бы $CT_a=I$, то
\[
  (CT_b)T_{\alpha_0-t}=I.
\]
Это дало бы левый обратный к $T_{\alpha_0-t}$ вопреки первой лемме.

\emph{2. Правый обратный.} Теперь выделим другой множитель:
\[
  a(t)=\left(1-\frac{\alpha_0}{t}\right)c(t),\qquad
  c(t)=\frac{q(t)}{t^{m-1}}.
\]
При $m=0$ последняя запись означает $c(t)=tq(t)$.
Символ $1-\alpha_0/t$ принадлежит $W_-$, поэтому соответствующий
оператор $A_0$ верхнетреуголен и $P_+A_0=P_+A_0P_+$.
Следовательно,
\[
  T_a=T_{1-\alpha_0/t}T_c.
\]
Если бы $T_aD=I$, то
\[
  T_{1-\alpha_0/t}(T_cD)=I.
\]
Это дало бы правый обратный к $T_{1-\alpha_0/t}$ вопреки второй лемме.
Обе возможности исключены.
\end{proof}

\begin{remark}
В первом случае линейный множитель поставлен справа, во втором ---
слева. Именно такое расположение позволяет вставить промежуточный
проектор. Для произвольных символов равенство $T_{bc}=T_bT_c$
не выполняется автоматически.
\end{remark}

\section{Необходимое условие для произвольного ядра}

\begin{theorem}
Пусть $a\in W$. Если $T_a$ обратим хотя бы с одной стороны,
то его символ нигде не равен нулю на единичной окружности.
\end{theorem}

\begin{proof}
Предположим противное: $a(\alpha_0)=0$ при $|\alpha_0|=1$.
По устойчивости односторонней обратимости существует $\delta>0$
такое, что все операторы $T$ с
\[
  \norm{T-T_a}<\delta
\]
имеют обратный с той же стороны, с которой он существует у $T_a$.
Если $T_a$ двусторонне обратим, можно использовать обычную
устойчивость обратимости. В рукописи подробно рассматривается
правый обратный; ниже оценка одинакова для обеих сторон.

Поскольку коэффициенты $a$ абсолютно суммируемы, выберем $N$ так, чтобы
\[
  \tau_N:=\sum_{|k|>N}|a_k|<\frac\delta2.
\]
Положим
\[
  a_N(t)=\sum_{k=-N}^{N}a_kt^k.
\]
Тогда $\norm{a-a_N}_W=\tau_N$ и
\[
  \norm{T_a-T_{a_N}}
  \leqslant\norm{A-A_N}=\norm{a-a_N}_W=\tau_N<\frac\delta2.
\]
Здесь использовано $\norm{P_+}=1$ и равенство нормы двусторонней
свёртки норме её ядра.

Однако $a_N$ не обязана иметь нуль в $\alpha_0$. Исправим только
постоянный коэффициент, определив
\[
  \boxed{b(t)=a_N(t)-a_N(\alpha_0)}.
\]
Это всё ещё многочлен Лорана, но теперь $b(\alpha_0)=0$ точно.
Так как $a(\alpha_0)=0$, получаем
\begin{align*}
  |a_N(\alpha_0)|
  &=|a_N(\alpha_0)-a(\alpha_0)|\\
  &\leqslant\max_{|t|=1}|a_N(t)-a(t)|
   \leqslant\norm{a_N-a}_W=\tau_N<\frac\delta2.
\end{align*}
На $X_+$ имеем $T_b=T_{a_N}-a_N(\alpha_0)I_{X_+}$. Поэтому
\begin{align*}
  \norm{T_b-T_a}
  &\leqslant\norm{T_{a_N}-T_a}+|a_N(\alpha_0)|\norm{I_{X_+}}\\
  &\leqslant\tau_N+|a_N(\alpha_0)|<\delta.
\end{align*}
По выбору $\delta$ оператор $T_b$ имеет односторонний обратный.
Но его ядро конечно, а $b(\alpha_0)=0$, что противоречит теореме
предыдущего раздела. Значит, исходное предположение неверно.
\end{proof}

\begin{remark}
Поправка $-a_N(\alpha_0)$ существенна: простое усечение ряда может
сдвинуть нуль с окружности. Поправка сохраняет нуль точно и при этом
увеличивает погрешность не более чем на величину хвоста ряда.
Для сжатых операторов используется оценка
$\norm{T_b-T_a}\leqslant\norm{b-a}_W$; равенство этих норм
для доказательства не требуется.
\end{remark}

\section{Полный критерий обратимости}

Объединяя достаточные и необходимые условия, получаем следующую
формулировку. Она собирает доказанные в рукописи результаты в одну таблицу.

\begin{theorem}
Для $a\in W$ оператор $T_a$ обладает односторонним обратным
тогда и только тогда, когда $a(t)\ne0$ на $\Gamma$.
Если символ ненулевой, тип обратимости определяется его индексом:
\[
  \boxed{T_a\text{ двусторонне обратим}
  \quad\Longleftrightarrow\quad
  a(t)\ne0\ (t\in\Gamma),\ \ind a=0}.
\]
\end{theorem}

\begin{center}
\renewcommand{\arraystretch}{1.35}
\begin{tabular}{|p{0.45\textwidth}|p{0.44\textwidth}|}
\hline
Условие на символ & Тип обратимости $T_a$\\\hline
$a(t)\ne0$ на $\Gamma$, $\ind a>0$ & Только слева\\\hline
$a(t)\ne0$ на $\Gamma$, $\ind a=0$ & Двусторонняя\\\hline
$a(t)\ne0$ на $\Gamma$, $\ind a<0$ & Только справа\\\hline
$a(\alpha_0)=0$ для некоторого $|\alpha_0|=1$ & Ни слева, ни справа\\\hline
\end{tabular}
\end{center}

Если у символа есть нуль на окружности, индекс в использованном
определении не задаётся. Поэтому последняя строка не является
случаем какого-либо значения индекса.

\subsection{Примеры применения критерия}

\emph{Дополнение для проверки знаков.} У функции $a(t)=t^2$
индекс $2$, поэтому $T_a=R^2$ обратим только слева.
У $a(t)=t^{-2}$ индекс $-2$, поэтому $T_a=L^2$ обратим только справа.

У $a(t)=2-t$ нуль расположен вне единичного круга, индекс равен
нулю. Оператор $2I-R$ обратим; например,
\[
  (2I-R)^{-1}=\frac12\sum_{n=0}^{\infty}\left(\frac R2\right)^n.
\]
Наконец, символ $a(t)=1-t$ обращается в нуль при $t=1$,
поэтому $I-R$ не имеет обратного ни с одной стороны.
Это последнее утверждение сильнее одной лишь необратимости.

\end{document}
